\section{Problématique et cadre théorique}

\begin{frame}{Problématique}
\begin{columns}[T]
    \column{.48\textwidth}
    \textcolor{navy}{\textbf{Cerveau humain}}
    \begin{itemize}
        \item $\sim$86 milliards de neurones
        \item $\sim$20 W
        \item Acquisition du langage en quelques années d'exposition
    \end{itemize}
    \column{.48\textwidth}
    \textcolor{navy}{\textbf{Modèle de langage}}
    \begin{itemize}
        \item $10^9$ à $10^{11}$ paramètres
        \item Entraînement sur $\sim 10^{12}$ mots
        \item Performance proche de l'humain en production
    \end{itemize}
\end{columns}
\vspace{1em}
\begin{block}{Question}
    Ces deux systèmes mettent-ils en œuvre des calculs comparables ?
\end{block}
\end{frame}

\begin{frame}{Précédents historiques}
\begin{itemize}
    \item 1943 : neurone formel (McCulloch \& Pitts) ; 1958 : perceptron (Rosenblatt)
    \item Réseaux convolutifs inspirés du cortex visuel (Hubel \& Wiesel)
    \item 2014 : des réseaux optimisés pour la reconnaissance d'objets prédisent l'activité du cortex visuel des primates
\end{itemize}
\vspace{0.5em}
\begin{block}{Extension au langage}
    Question abordée depuis environ 2018 avec les modèles de langage.
\end{block}
\src{Yamins et al., PNAS 2014}
\end{frame}

\begin{frame}{Les trois niveaux de Marr}
\begin{columns}[T]
    \column{.52\textwidth}
    \begin{itemize}
        \item<1-> \textbf{Computationnel} : \emph{quel problème} est résolu ?\\
              {\small ex. : une addition}
        \item<2-> \textbf{Algorithmique} : \emph{quelles représentations et quelles procédures} ?\\
              {\small ex. : chiffres arabes et retenue}
        \item<3-> \textbf{Implémentationnel} : \emph{substrat} physique\\
              {\small ex. : neurones, transistors}
    \end{itemize}
    \column{.46\textwidth}
    \centering
    \resizebox{\linewidth}{!}{%
    \begin{tikzpicture}[
        n/.style={circle, draw=navy, fill=navy!15, minimum size=0.42cm, inner sep=0pt},
        e/.style={navy, thick},
        lab/.style={font=\small\bfseries, text=navy, anchor=west}]
        \node<1->[n] (c) at (0,0) {};
        \node<1->[lab] at (2.4,0) {Computationnel};
        \node<2->[n] (a1) at (-1.2,-1.4) {};
        \node<2->[n] (a2) at (0,-1.4) {};
        \node<2->[n] (a3) at (1.2,-1.4) {};
        \node<2->[lab] at (2.4,-1.4) {Algorithmique};
        \draw<2->[e] (c)--(a1); \draw<2->[e] (c)--(a2); \draw<2->[e] (c)--(a3);
        \foreach \x/\i in {-1.8/1,-0.6/2,0.6/3,1.8/4}{
            \node<3->[n] (i\i) at (\x,-2.8) {};
            \draw<3->[e] (a2)--(i\i);}
        \node<3->[lab] at (2.4,-2.8) {Implémentation};
        \node<3->[font=\small\bfseries, text=navy] at (0,-3.5) {Relation un-à-plusieurs};
    \end{tikzpicture}}
\end{columns}
\end{frame}

\begin{frame}{Marr : humains et modèles de langage}
\centering
\resizebox{\textwidth}{!}{%
\begin{tikzpicture}[
    n/.style={circle, draw=navy, fill=navy!15, minimum size=0.42cm, inner sep=0pt},
    e/.style={navy, thick},
    lab/.style={font=\small\bfseries, text=navy, anchor=west}]
    \node[n] (c) at (0,0) {};
    \foreach \x/\i in {-4/1,-2/2,0/3,2/4,4/5}{\node[n] (a\i) at (\x,-1.5) {};}
    \foreach \i in {2,3,4}{\draw[e] (c)--(a\i);}
    \draw[lscp, ultra thick] (c)--(a1);
    \draw[blue2, ultra thick] (c)--(a5);
    \foreach \x/\i in {-5.2/1,-4.4/2,-3.6/3,-2.8/4}{\node[n] (m\i) at (\x,-3.2) {}; \draw[e] (a1)--(m\i);}
    \foreach \x/\i in {2.8/1,3.6/2,4.4/3,5.2/4}{\node[n] (h\i) at (\x,-3.2) {}; \draw[e] (a5)--(h\i);}
    \node[lab, text=lscp!70!black] at (-5.2,-3.9) {Modèles de langage};
    \node[lab, text=blue2] at (2.6,-3.9) {Cerveaux humains};
    \node[lab] at (6.2,0) {Computationnel};
    \node[lab] at (6.2,-1.5) {Algorithmique};
    \node[lab] at (6.2,-3.2) {Implémentation};
\end{tikzpicture}}

\vspace{0.4em}
\textbf{Hypothèse de divergence} : implémentations très différentes, éventuellement objectifs différents ; les algorithmes diffèrent.
\end{frame}

\begin{frame}{Marr : humains et modèles de langage}
\centering
\resizebox{\textwidth}{!}{%
\begin{tikzpicture}[
    n/.style={circle, draw=navy, fill=navy!15, minimum size=0.42cm, inner sep=0pt},
    e/.style={navy, thick},
    lab/.style={font=\small\bfseries, text=navy, anchor=west}]
    \node[n] (c) at (0,0) {};
    \foreach \x/\i in {-4/1,-2/2,0/3,2/4,4/5}{\node[n] (a\i) at (\x,-1.5) {};}
    \foreach \i in {1,2,4,5}{\draw[e] (c)--(a\i);}
    \draw[lscp, ultra thick] (c) to[bend right=28] (a3);
    \draw[blue2, ultra thick] (c) to[bend left=28] (a3);
    \foreach \x/\i in {-5.2/1,-4.4/2,-3.6/3,-2.8/4}{\node[n] (m\i) at (\x,-3.2) {}; \draw[e] (a3)--(m\i);}
    \foreach \x/\i in {2.8/1,3.6/2,4.4/3,5.2/4}{\node[n] (h\i) at (\x,-3.2) {}; \draw[e] (a3)--(h\i);}
    \node[lab, text=lscp!70!black] at (-5.2,-3.9) {Modèles de langage};
    \node[lab, text=blue2] at (2.6,-3.9) {Cerveaux humains};
    \node[lab] at (6.2,0) {Computationnel};
    \node[lab] at (6.2,-1.5) {Algorithmique};
    \node[lab] at (6.2,-3.2) {Implémentation};
\end{tikzpicture}}

\vspace{0.4em}
\textbf{Hypothèse de convergence} : un \emph{même} algorithme, implémenté sur deux substrats différents.
\end{frame}

\begin{frame}{Inférer le niveau algorithmique}
\begin{itemize}
    \item L'algorithme n'est pas directement observable : il est \textbf{inféré}
    \item \textbf{Comportement} : temps de lecture (lecture auto-rythmée, oculométrie)
    \item \textbf{Activité cérébrale} : motifs d'activité reflétant les représentations utilisées
    \item On compare donc modèles et humains \textbf{au niveau algorithmique}
\end{itemize}
\vspace{0.4em}
\begin{block}{Intérêt des modèles d'IA}
    Sous l'hypothèse de convergence, un modèle est une hypothèse explicite sur l'algorithme cérébral, entièrement observable et manipulable (ablation, intervention), ce que le cerveau permet peu.
\end{block}
\src{Marr, Vision, 1982}
\end{frame}

\begin{frame}{Domaines abordés}
\begin{columns}[T]
    \column{.48\textwidth}
    \begin{block}{Modèle du système cognitif}
        \small
        Tester des hypothèses sur le calcul cérébral\\[0.3em]
        Vision et récompense ; langage (alignement) ; comportement ; mécanismes
    \end{block}
    \column{.48\textwidth}
    \begin{block}{Source d'inspiration}
        \small
        Le cerveau guide la conception des modèles\\[0.3em]
        Signal cérébral comme supervision ; apprentissage à échelle développementale
    \end{block}
\end{columns}
\vspace{0.8em}
Pour chaque domaine : \textbf{potentiel} et \textbf{état actuel} des résultats.
\end{frame}
